\begin{phasebox}
\textbf{هدف:} آماده‌سازی یک محیط محلی قابل تکرار.\\
\textbf{پیش‌نیاز:} دسترسی به برد و ابزار نصب سیستم‌عامل.\\
\textbf{معیار پایان:} نسخه ابزارها ثبت و تست مستقل هر ابزار موفق باشد.
\end{phasebox}

\section{ابزارهای پایه}
ابزارهای اولیه پروژه عبارت‌اند از:

\begin{itemize}
    \item
    \lr{Python}
    و محیط مجازی آن؛
    \item
    \lr{SQLite}
    برای ذخیره محلی؛
    \item
    \lr{Mosquitto}
    و ابزارهای خط فرمان آن؛
    \item
    \lr{Arduino IDE}
    برای ساخت و ارسال برنامه برد؛
    \item
    \lr{Git}
    برای ثبت تغییرات؛
    \item
    \lr{XeLaTeX}
    و
    \lr{TeXstudio}
    برای گزارش.
\end{itemize}

\section{آزمایش مستقل}
\begin{itemize}
    \item اجرای یک فایل کوچک پایتون در محیط مجازی؛
    \item ساخت یک پایگاه داده موقت و خواندن یک مقدار از آن؛
    \item ارسال و دریافت یک پیام محلی با ابزارهای
    \lr{Mosquitto}
    \item ساخت و ارسال یک برنامه خروجی سریال برای برد؛
    \item تولید فایل
    \lr{\texttt{main.pdf}}
    در کنار
    \lr{\texttt{main.tex}}
\end{itemize}


\subsection{راه‌اندازی آردویینو}
برای استفاده از کد آردویینو دو راه اساسی را می‌توان استفاده کرد. راه اول و راه ساده‌تر استفاده از خود 
\lr{Arduino IDE}
است که به سادگی می‌توان پروژه را ساخت و ازش استفاده کرد.
راه دیگر استفاده از 
\lr{PlatformIO}
است که به کمک آن قادر خواهیم بود که پروژ خود را داخل 
\lr{VSCode}
پیاده سازی کنیم و کدها را داخل آن بنویسی و به سادگی می‌توانیم به مدیریت فایل‌ها بپردازیم.
\section{خروجی ثبت‌شونده}
نسخه دقیق ابزارها، فرمان نصب تأییدشده و نتیجه هر آزمایش در این فصل ثبت می‌شود.
